In this section we empirically test the implications of our model.
The predictions we take to the data are those of the short-run problem of Section~\ref{sec:shortrun}, in which the boundaries of jobs are fixed and AI acts only on how the steps within them are executed.
This is a deliberate restriction rather than a limitation of the theory.
The long-run margins of Section~\ref{sec:longrun} require firms to reorganize and workers to retrain, and not enough time has passed since AI tools became widely available for these adjustments to appear in the data we observe.
What we can observe is where AI is being deployed inside workflows as they are currently arranged, which is precisely what the short-run problem speaks to.
The three predictions we test are: 
(1) AI-executed steps tend to appear next to each other in the production sequence, forming AI chains; 
(2) a step is more likely to be executed by AI when its neighbors are AI-executed; and 
(3) workflows in which AI-exposed steps are more dispersed have fewer of those steps actually executed by AI.

Our analysis requires three attributes of each step: its AI exposure status (exposed or unexposed to AI), which is how we measure whether a step is AI-able in the sense of Section~\ref{sec:results}, its realized mode of execution (manual or AI-executed), and its position in the production sequence.
We draw on five sources:
\begin{enumerate}[itemsep=0pt, parsep=0pt]
    \item The O*NET 27.3 Database \citep{onet27_3},
    \item Human-generated AI exposure labels for O*NET tasks from \cite{eloundou2023gpts},
    \item Realized AI execution labels for O*NET tasks from the Anthropic Economic Index \citep{handa2025economictasksperformedai},
    \item A GPT-generated task ordering for each O*NET occupation,
    \item The Process Classification Framework (PCF) of the American Productivity \& Quality Center \citep{apqc_pcf}.
\end{enumerate}
The first four build our O*NET sample, in which a workflow is an occupation, and the fifth builds our PCF sample, in which a workflow is a process group.
All sources are existing (and publicly available) datasets apart from the O*NET task ordering, which we generate ourselves.

The O*NET database is the standard decomposition of work in the United States.
It breaks each of more than 850 occupations into a list of constituent tasks, roughly 18,000 in total, and has served as the backbone of a large empirical literature on the task content of jobs and their exposure to new technologies \citep[e.g.,][]{autor2013putting, deming2017growing, brynjolfsson2018can, webb2020impact, felten2021occupational}.
\cite{eloundou2023gpts} map these O*NET tasks to AI exposure by developing a rubric asking whether access to an LLM would reduce the time required to perform a task by at least half while preserving output quality, and applying it to every task using both human annotators and GPT-4.
We use their human-generated labels and treat tasks in their E1 and E2 categories as exposed to AI and all other tasks as unexposed.
Whereas these labels capture what AI could in principle do, the Anthropic Economic Index \citep{handa2025economictasksperformedai} captures what AI actually does.
It maps millions of conversations between users and Claude, Anthropic's AI chatbot, to O*NET tasks and classifies each conversation by whether the interaction augments or automates the task, providing a task-level measure of realized AI usage.\footnote{One feature of the data shapes how we use it throughout. The model's automated versus augmented distinction is defined by production sequencing, whereas the Anthropic labels are assigned on task-level feasibility and are agnostic to workflow position, so the data occasionally shows an automated task followed by a manual one, which the model's definition of a chain does not allow. We therefore do not distinguish the two labels in what follows and treat both as indicating AI execution. Appendix~\ref{app:data_prep} documents the discrepancy and provides examples.}

The PCF catalogues business processes, with APQC's member working groups recording the activities that compose each process in the order they are performed.
Pooling APQC's Cross-Industry framework with its industry-specific ones leaves $13{,}482$ steps in $525$ process groups.
Since AI exposure and execution are defined on O*NET tasks, we carry those labels onto PCF steps by matching each step to its nearest O*NET task using natural language embeddings, keeping unmatched steps in sequence and coding them as neither AI-exposed nor AI-executed.
The label transfer is lossy, leaving $12\%$ ($4\%$) of PCF steps AI-exposed (AI-executed) against $44\%$ ($13\%$) in the O*NET sample.

The two corpora complement each other in our analysis.
The PCF's workflow orderings are recorded by practitioners and thus native to the dataset.
However, the PCF covers only operational business processes, and its AI exposure and execution content is thin because of the lossy PCF--O*NET match.
O*NET, on the other hand, has broader coverage of economic activity but lists an occupation's tasks without the order in which they are performed.
To make the O*NET dataset suitable for our analysis, we construct workflow orderings for it by prompting GPT-5-mini \citep[through Expected Parrot;][]{Horton2024EDSL} with each occupation's complete set of tasks and asking for the most reasonable order in which they would be performed in practice.
Appendices~\ref{app:data_prep} and~\ref{app:external_validation} respectively give the construction details of the O*NET and PCF samples.

Because each of our tests is defined over adjacent steps, the O*NET results, unlike the PCF results, rest on LLM-generated orderings, and we probe them in three robustness exercises.
First, we show that the task orderings, and by extension our conclusions, change little across ten deliberately varied alternative prompts (Appendix~\ref{app:gptPrompts_robustness}).
Second, we show that, applied to the PCF and to eight process-mining event logs \citep{roadfines_log, sepsis_log, bpi2017_log, bpi2020_log} whose orderings are known, our sequencing procedure restores the correct workflow order to a substantial degree against ground truth (Appendix~\ref{app:external_validation}).
Third, restricting the O*NET sample to progressively more frequently performed tasks, which drops rarely performed ``filler'' tasks that can separate steps adjacent in practice, leaves our conclusions unchanged (Appendix~\ref{app:frequency_robustness}).
Our O*NET results are therefore unlikely to be artifacts of the prompt, the sequencing instrument, or the way O*NET catalogues tasks.

The structure of our theoretical framework closely aligns with the structure of both datasets.
Each production step in the model corresponds to a PCF step or an O*NET task, and each workflow corresponds to the model's step sequence.
In the absence of AI, a step (that is, a PCF step or an O*NET task) is executed manually and maps directly into a task under our model's definition.
Thus, without AI, a PCF step or an O*NET task can be viewed as both a step and a task, which matches how both datasets list their units of work.
Under AI execution, however, multiple steps, corresponding to a subset of PCF steps or O*NET tasks, may be chained together and executed jointly, forming a composite task as defined in Section~\ref{sec:shortrun}.
In the rest of this section, we use the terms step, PCF step, O*NET task, and task interchangeably and expect the reader to keep these distinctions in mind.
We similarly use the term workflow for both an O*NET occupation and a PCF process group without restating the mapping each time.


\subsection{Prediction \#1: Tendency to Have Runs of Consecutive AI-executed Tasks}
\label{sec:chainLength_prediction}

The first prediction of the model that we test is that AI-executed steps appear in contiguous blocks rather than being randomly scattered in the workflow.
This prediction follows directly from Definition~\ref{def:ai_chain} on how AI chains are formed in the model.
If firms deploy AI, chaining can lower costs because the worker verifies the output only once, at the end of the chain, rather than after every step.

We begin with the O*NET sample.
We compute the observed average AI chain length and benchmark it against two randomization placebos.
First, we permute task positions within each occupation's workflow.
Second, we reassign tasks, each carrying its AI-execution label, randomly across occupations within the same major SOC group, preserving each occupation's task count and randomizing task positions within each occupation.
The two placebos isolate different margins: the first holds each occupation's task composition fixed and randomizes the ordering of its workflow, while the second holds each occupation's size and each major group's AI intensity fixed and randomizes which tasks, and hence which execution labels, an occupation contains.

Panel~(a) of Figure~\ref{fig:aiChains_graphs_def1} overlays the two placebo distributions of average AI chain length, one across 1{,}000 task position permutations and one across 1{,}000 task AI-execution label reassignment permutations, and compares both with the observed mean AI chain length, marked by the dashed vertical line.
We observe that the average AI chain length in the data is 1.45 and is noticeably larger than in both placebo distributions.
This indicates that AI-executed tasks tend to cluster within occupations, forming longer AI chains than would arise under either type of randomization.
At the same time, the modest magnitude of 1.45 indicates that long AI chains remain rare, perhaps reflecting the early stage of AI adoption; our model predicts that chains will lengthen as AI quality improves.
\begin{figure}[!t]
\caption{Average Length of AI Chains: Actual versus Placebo Datasets}
\label{fig:aiChains_graphs_def1}

\centering
\begin{subfigure}[b]{0.49\textwidth}
    \centering
    \subcaption{O*NET Occupations}
    \label{fig:aiChains_onet}
    \includegraphics[width=\textwidth]{plots/chainLength_overlay/aiChains_chainLength_overlay_onet.png}
\end{subfigure}
\hfill
\begin{subfigure}[b]{0.49\textwidth}
    \centering
    \subcaption{PCF Process Groups}
    \label{fig:aiChains_pcf}
    \includegraphics[width=\textwidth]{plots/chainLength_overlay/aiChains_chainLength_overlay_apqc.png}
\end{subfigure}

\par\vspace{\fignotesskip}
\begin{minipage}{1\linewidth}
\begin{spacing}{0.2}
{\fontsize{9.5pt}{10pt}\selectfont Notes: Red dashed lines show the observed average AI chain length in the actual dataset.
Panel~(a) uses the O*NET sample, in which a workflow is an occupation and the ordering is the GPT-imputed task sequence; Panel~(b) uses the PCF sample, in which a workflow is a process group and the ordering is the one documented by practitioners, with exposure and execution labels transferred from each step's nearest O*NET task.
Each panel overlays both placebo distributions, each over $1{,}000$ draws.
The orange histogram randomizes the positions of steps within each workflow, and the blue histogram reassigns steps, each carrying its execution label, across workflows within the same group, preserving workflow sizes; the grouping is the major SOC group in Panel~(a) and the PCF Category in Panel~(b).
The horizontal scale differs between the panels, since only $4\%$ of PCF steps carry an AI-execution label after the transfer against $13\%$ in O*NET, which lowers the attainable chain length throughout the PCF sample.}
\end{spacing}
\end{minipage}
\end{figure}

Next we repeat the exercise on the PCF sample, substituting its own units throughout, with a process group standing in for an occupation, the activity sequence practitioners documented for the GPT-imputed task sequence, and the PCF Category for the major SOC group within which the second placebo reassigns.
Nothing else about the test changes, and Panel~(b) of Figure~\ref{fig:aiChains_graphs_def1} reports the same verdict.
The average AI chain there spans $1.16$ steps, against $1.45$ in O*NET.
The smaller magnitude is mechanical, since AI execution in the PCF is thinner than in O*NET, which leaves far less opportunity for long runs to form and puts the ceiling on the statistic well below the O*NET sample's.
What the PCF sample can still ask is whether the runs that do form are longer than the same labels arranged at random along the same sequences, and they are.
Because no language model placed these steps in order, the contiguity here is a property of the documented workflow and not of our sequencing procedure.



\subsection{Prediction \#2: AI Execution Is More Likely Next to AI-executed Steps}
\label{sec:DWA_prediction}

The second prediction of the model that we test is that when the same step appears in two occupations, in one sitting between strong AI performers and in the other located between manual ones, it is more likely to be executed by AI in the first occupation.
This is the empirical content of Proposition~\ref{prop:ca_local}.
Intuitively, when a step is surrounded by AI-executed steps, there are local gains from grouping it with its neighbors as part of a long AI chain that are absent when the step is situated next to manual steps.
For example, step X is more likely to be executed by AI in the first occupation below than in the second.

\vspace{0.6em}
\begin{adjustbox}{center}
\begin{tikzpicture}[scale=0.93, transform shape,
    node distance=1cm,
    every node/.style={rectangle, rounded corners, draw, minimum width=1.9cm, minimum height=0.9cm, align=center},
    human/.style={fill=olive!20},
    ai/.style={fill=blue!20},
    xnode/.style={fill=teal!20},
    >=latex
]

% ======================
% OCCUPATION 1
% ======================

\node[draw=none, anchor=west] (L1) at (0,0) {\textbf{Occupation 1:}};

% Place chain on its own clean baseline
\begin{scope}[xshift=3.5cm, yshift=0cm]

\node[draw=none] (Dots1L) at (0,0) {\ \ $\cdots \cdots$};
\node[ai, right=1 of Dots1L] (A1a) {AI};
\node[xnode, right=1 of A1a] (X1) {X};
\node[ai, right=1 of X1] (A1b) {AI};
\node[draw=none, right=1 of A1b] (Dots1R) {$\cdots \cdots$};

\draw[->, thick] (Dots1L.east) -- (A1a);
\draw[->, thick] (A1a) -- (X1);
\draw[->, thick] (X1) -- (A1b);
\draw[->, thick] (A1b) -- (Dots1R.west);

\end{scope}


% ======================
% OCCUPATION 2
% ======================

\node[draw=none, anchor=west] (L2) at (0,-1.4) {\textbf{Occupation 2:}};

% Place chain on its own line (same centering)
\begin{scope}[xshift=3.5cm, yshift=-1.4cm]

\node[draw=none] (Dots2L) at (0,0) {\ \ $\cdots \cdots$};
\node[human, right=1 of Dots2L] (H2a) {Manual};
\node[xnode, right=1 of H2a] (X2) {X};
\node[human, right=1 of X2] (H2b) {Manual};
\node[draw=none, right=1 of H2b] (Dots2R) {$\cdots \cdots$};

\draw[->, thick] (Dots2L.east) -- (H2a);
\draw[->, thick] (H2a) -- (X2);
\draw[->, thick] (X2) -- (H2b);
\draw[->, thick] (H2b) -- (Dots2R.west);

\end{scope}

\end{tikzpicture}
\end{adjustbox}
\vspace{-1.4em}

Testing this prediction requires identifying steps that appear across different workflows.
Because O*NET tasks are occupation-specific, no single task is shared across occupations in the data.
To work around this, we rely on O*NET's broader categories of Detailed Work Activities (DWAs), which group conceptually similar tasks across occupations.
In our baseline O*NET sample, we drop tasks that are mapped to multiple DWAs.
We further restrict attention to DWAs that contain tasks appearing in more than one occupation.
Together, these restrictions reduce the sample to 10{,}708 tasks spread across 1{,}748 DWAs.
We treat the remaining tasks as instances of the same step appearing across different occupations and estimate the following regression\footnote{In a robustness exercise we refine our definition of ``similar'' tasks.
The resulting estimates closely mirror the patterns observed in the full sample in both sign and magnitude.
See Appendix~\ref{app:additional_robustness} for details.}:
% Broken after the \beta_2 term: on one line the regression is 523pt wide against a
% 470pt measure and the closing \Big) and comma ran past the right margin. \small is
% not enough on its own (477pt), and the two rows come to 312pt and 205pt at full size.
% \Big( and \Big) are explicit-size delimiters, not \left ... \right, so they can sit on
% different rows of the alignment.
\begin{equation}
\begin{aligned}
    \Pr\!\left( \text{is\_ai}_{k} = 1 \mid X_k \right)
    ={}
    &\Lambda\!\Big(
        \beta_0
        + \beta_1 \, \text{prev2\_is\_ai}_{k}
        + \beta_2 \, \text{prev\_is\_ai}_{k} \\
    &\qquad
        + \beta_3 \, \text{next\_is\_ai}_{k}
        + \beta_4 \, \text{next2\_is\_ai}_{k}
    \Big),
\end{aligned}
    \label{eq:DWA_regression_ai}
\end{equation}
where $\Lambda$ is the logistic CDF and \emph{$\text{is\_ai}_{k}$} indicates whether step $k$ is executed by AI.
The variables \emph{$\text{prev2\_is\_ai}_{k}$} and \emph{$\text{prev\_is\_ai}_{k}$} equal 1 if the step two positions before or immediately before step $k$ is AI-executed, respectively, and \emph{$\text{next\_is\_ai}_{k}$} and \emph{$\text{next2\_is\_ai}_{k}$} are defined analogously for the steps immediately after and two positions after step $k$.
In the implementation we also control for step $k$'s AI exposure status and the number of steps in its occupation.

Table~\ref{tab:DWA_regression_aiExecution_mainSample} reports average marginal effects from estimating Equation~\eqref{eq:DWA_regression_ai} for a series of specifications.
\begin{table}[!t]
\centering
\caption{Neighboring Tasks' AI Execution Status and a Task's AI Execution}
\label{tab:DWA_regression_aiExecution_mainSample}
    \footnotesize
\resizebox{\textwidth}{!}{
% --- LaTeX Table for full_0 ---
\begin{tabular}{lcccccc}
\toprule
 & \multicolumn{6}{c}{Probability that Focal Task ($k$) is AI-executed} \\
\cmidrule(lr){2-7}
 & (1) & (2) & (3) & (4) & (5) & (6) \\
\midrule
Task ($k-2$) is AI-executed & 0.07*** & 0.01 & 0.00 & -0.01 & 0.06* & -0.01 \\
 & (0.02) & (0.01) & (0.01) & (0.03) & (0.03) & (0.03) \\
\addlinespace
Task ($k-1$) is AI-executed & 0.12*** & 0.06*** & 0.05*** & 0.05** & 0.13*** & 0.05** \\
 & (0.02) & (0.01) & (0.01) & (0.02) & (0.03) & (0.02) \\
\addlinespace
Task ($k+1$) is AI-executed & 0.12*** & 0.06*** & 0.05*** & 0.04** & 0.10*** & 0.04** \\
 & (0.02) & (0.01) & (0.01) & (0.02) & (0.02) & (0.02) \\
\addlinespace
Task ($k+2$) is AI-executed & 0.05*** & 0.00 & -0.01 & 0.00 & 0.04 & 0.00 \\
 & (0.02) & (0.01) & (0.01) & (0.02) & (0.03) & (0.02) \\
\addlinespace
\midrule
Pseudo $R^2$ & 0.112 & 0.173 & 0.171 & 0.196 & 0.030 & 0.197 \\
Observations & 10,708 & 10,708 & 9,861 & 4,096 & 4,096 & 4,096 \\
SOC Group FE &  & Major & Minor &  &  &  \\
DWA FE &  &  &  & \checkmark &  & \checkmark \\
NumTasks in DWA-Occupation Control &  &  &  &  & \checkmark & \checkmark \\
\bottomrule
\end{tabular}

}
\begin{minipage}{1\linewidth}
\begin{spacing}{0.2}
{\fontsize{9.5pt}{10pt}\selectfont Notes: Clustered standard errors in parentheses. *** p$<$0.01, ** p$<$0.05, * p$<$0.1.
This table reports average marginal effects from estimating Equation~\eqref{eq:DWA_regression_ai}.
The dependent variable in all specifications is an indicator for whether task $k$ is AI-executed ($\text{is\_ai}_{k}$).
The sample is restricted to similar tasks across occupations, identified as tasks belonging to the same DWA in the O*NET dataset.
All specifications control for the AI exposure status of task $k$ and for the total number of tasks in task $k$'s occupation.
Standard errors are bootstrapped using $B=200$ replications and clustered at the DWA level.
The estimating sample varies across columns. Minor-group fixed effects in column~(3) drop groups with no within-group variation in the outcome, removing 847 tasks. The DWA fixed effects in columns~(4) and~(6) are identified only from DWAs containing both an AI-executed and a non-AI-executed task, a condition met by 534 of the 1{,}748 DWAs in the sample, leaving 4{,}096 tasks; column~(5) carries no DWA fixed effects but is estimated on that same restricted sample so that it is directly comparable with column~(4). The number of DWA clusters is 1{,}748 in columns~(1)--(2), 1{,}705 in column~(3), and 534 in columns~(4)--(6).}
\end{spacing}
\end{minipage}\end{table}
The baseline in column~(1) shows that AI execution anywhere in the focal task's local context is associated with a higher probability that the task is AI-executed, at one and at two positions away alike, with the effects of the adjacent steps roughly twice the size of those two positions out.
The remaining columns ask how much of this survives more demanding comparisons.
Columns~(2)--(3) absorb occupation family differences with SOC major and minor group fixed effects, and column~(4) restricts the variation to within detailed work activities with DWA fixed effects, addressing differences in AI-ability across types of task.
Columns~(5)--(6) repeat the exercise controlling for the number of same-DWA tasks in the focal task's occupation, which rules out mechanical proximity inflating the estimates when several same-DWA tasks sit in one workflow; columns~(4)--(6) are all estimated on the same restricted sample, so they are comparable with one another rather than with column~(1).
Across the specifications that absorb occupation-family or DWA heterogeneity, columns~(2)--(4) and~(6), the pattern is the same.
The immediate neighbor effect attenuates to between $0.04$ and $0.06$, relative to the baseline for columns~(2)--(3) and to the restricted-sample estimate of column~(5) for columns~(4) and~(6), but stays statistically and economically meaningful, while the effects two positions away fall to between $-0.01$ and $0.01$ and lose significance.
Column~(5) carries no fixed effects and, estimated on the restricted sample of column~(4), reproduces the baseline pattern rather than the attenuated one.
That the baseline estimates attenuate tells us they partly reflect systematic differences across occupation families; that the immediate neighbor effect persists tells us the local pattern is not merely such a difference.

Because we have no external benchmark for what constitutes a large effect here, the magnitudes are hard to read on their own.
We therefore re-estimate Equation~\eqref{eq:DWA_regression_ai} on 1,000 placebo datasets that randomize task positions within occupations, for the four specifications in columns~(1)--(4) of Table~\ref{tab:DWA_regression_aiExecution_mainSample}, and report the result in Figure~\ref{fig:DWA_regression_aiExecution_mainSample} of Appendix~\ref{app:tables_and_figures}.
Randomizing positions strips the ordering of any information, so the placebo distributions show what these coefficients would look like when the position of a task in the workflow does not matter.
In each case, the actual orderings deliver stronger immediate-neighbor effects and weaker distant-neighbor effects than the placebos.
This implies that local work context and proximity do work that the reshuffled orderings cannot reproduce by chance.

Overall, the results indicate that AI execution of immediate neighbors of the focal task is associated with a higher likelihood that the task is executed by AI, whereas the execution status of more distant neighbors matters much less, especially after controlling for a range of relevant factors.
This aligns with the relatively short average AI chain length observed in the data (mean length 1.45 reported in Subsection~\ref{sec:chainLength_prediction}), implying that tasks two positions away rarely fall in the same chain as the focal task.
It is possible that with improvements in AI quality, chains become longer and execution status of more distant neighbors becomes more relevant for the execution mode of the focal task.

For this prediction, we do not attempt a counterpart on the PCF sample, where too few steps are observed with both execution outcomes to support the comparison. See Appendix~\ref{app:external_validation} for details.



\subsection{Prediction \#3: AI Execution Is Lower Where AI-able Steps Are Dispersed}
\label{sec:fragmentation_prediction}

The two tests so far were local, in that they operated over a small window around a step.
The last prediction of the model that we test is how the positioning of steps over the whole workflow matters for AI execution.
Specifically, we test whether realized AI execution in a production sequence varies with how clustered or dispersed its AI-exposed steps are, and not merely with how many of its steps are exposed.
This is the empirical content of Proposition~\ref{prop:fragmentation}.

To measure the fragmentation of AI-exposed steps in the data, we adapt the fragmentation index of Section~\ref{sec:fragmentation} to our empirical setting.
The \emph{empirical fragmentation index} (EFI) is the ratio of the number of \emph{potential tasks}, as defined by our theoretical model, to the total number of steps, which is fixed and is measured by the raw count of steps in the workflow.
Consider a workflow with five steps.
If none of the five steps are exposed to AI, then the number of tasks in the theoretical sense is five, and the EFI equals $5/5 = 1$.
If instead two consecutive steps are exposed to AI (and thus can form an AI chain of length two when all AI-exposed steps turn into AI-executed steps), the number of tasks in the model falls to four and the EFI becomes $4/5 = 0.8$.
Thus, the empirical fragmentation index decreases as more AI-exposed steps can potentially be consolidated into longer AI chains.\footnote{The EFI is, after normalizing by the number of steps, the special case of the realized fragmentation expressed in \eqref{eq:fragmentation} in which every step takes one unit of time in either mode, $\manualTime{i} = \AItime{i} = 1$, and AI either always or never succeeds at a step, $q_i \in \{0,1\}$.
It therefore captures the arrangement of a workflow's AI-able steps, which is what we observe in the data, and not the differences in time cost and reliability that the theoretical index also prices.
Since $\manualTime{i} = 1$ satisfies the hypothesis of \eqref{eq:fragmentation_closed_form}, the EFI is exactly one minus the share of adjacent AI-able pairs in the workflow, so it inherits the ordering that expression establishes.}

Two components govern the empirical fragmentation index, and only one is the subject of Prediction~\#3.
Write $m_w$ for the number of steps in a workflow, $k_w$ for how many of them are AI-exposed, and $r_w$ for the number of maximal blocks of consecutive AI-exposed steps they form, each block being one potential AI chain.
EFI decomposes exactly as
\begin{equation}
    \text{EFI}_w \;=\; \frac{m_w - (k_w - r_w)}{m_w} \;=\; 1 - \underbrace{\frac{k_w}{m_w}}_{\text{exposure level}} \;+\; \underbrace{\frac{r_w}{m_w}}_{\text{arrangement}},
    \label{eq:efi_decomposition}
\end{equation}
where the last term is the object of interest.
Holding $k_w$ fixed, $r_w$ runs from $1$ when the AI-exposed steps are perfectly clustered to $k_w$ when no two of them are adjacent.
To read the index as a measure of dispersion, we must control for the share of AI-exposed steps, which nets out the level term and leaves the arrangement component.
We therefore estimate the following regression:
\begin{equation}
\begin{aligned}
    \text{ai\_execution}_{w} = {} & \beta_0 + \beta_1 \, \text{ai\_exposure}_{w} + \beta_2 \, \text{empirical\_fragmentation\_index}_{w} \\
    & + \beta_3 \, \text{num\_ai\_able\_steps}_{w} + \varepsilon_w,
\end{aligned}
    \label{eq:fragmentation_index_regression}
\end{equation}
where \emph{$\text{ai\_execution}_{w}$} is the share of steps executed by AI in workflow $w$, \emph{$\text{ai\_exposure}_{w}$} is the share of AI-exposed steps, \emph{$\text{empirical\_fragmentation\_index}_{w}$} is the index defined above, and \emph{$\text{num\_ai\_able\_steps}_{w}$} is the number of AI-able steps in workflow $w$.\footnote{Controlling for workflow length $m_w$ instead, or dropping the count control altogether, moves the fragmentation coefficient by at most $0.04$ standard deviations and changes no verdict.}
The model predicts that workflows with higher AI exposure should have a higher share of steps executed by AI ($\beta_1 > 0$), and that, conditional on exposure, AI execution should be lower where AI-able steps are more dispersed ($\beta_2 < 0$).

Table~\ref{tab:fragmentation_index_regression_exposure} reports the results.
\begin{table}[!t]
    \caption{AI-exposed Step Fragmentation and AI Execution Outcomes}
    \label{tab:fragmentation_index_regression_exposure}
    \footnotesize
    \centering
    \adjustbox{max width=\textwidth}{\input{tables/fragmentation_index_regression_combined.tex}}
    \begin{minipage}{1\linewidth}
\begin{spacing}{0.2}
{\fontsize{9.5pt}{10pt}\selectfont Notes: Standardized coefficients. Heteroskedasticity-robust standard errors in parentheses. *** p$<$0.01, ** p$<$0.05, * p$<$0.1.
This table reports results from estimating Equation~\eqref{eq:fragmentation_index_regression}.
The dependent variable in all specifications is the share of steps executed by AI in the workflow.
``Share of AI-exposed Steps'' is the share of steps carrying an E1 or E2 label, the same label set from which the Empirical Fragmentation Index is constructed; the index itself captures how dispersed AI-exposed steps are across the workflow.
All specifications control for the number of AI-able steps in the workflow. 
Columns~(1)--(3) use the O*NET sample, where a workflow is an occupation and the ordering is the GPT-imputed task sequence; columns~(4)--(6) use the PCF sample, where a workflow is a process group and the ordering is the one documented by practitioners.
In the PCF sample, exposure and execution labels are transferred from each step's nearest O*NET task.
All variables are standardized to unit standard deviation within their own sample, so the reported coefficients are in standard-deviation units and are not directly comparable in levels across the two samples.}
\end{spacing}
\end{minipage}\end{table}
In each sample we estimate three specifications: no fixed effects, and then two groupings of workflows (SOC major and minor occupation groups in O*NET, PCF Category and Framework in the PCF).
The fixed effects ensure the results are not driven by a narrow subset of occupation or process families.
The exposure coefficient behaves as the model predicts in every column of both samples.
A one-standard-deviation increase in the share of AI-exposed steps is associated with a $0.39$ to $0.49$ standard-deviation increase in the AI-executed share in O*NET and a $0.28$ to $0.42$ standard-deviation increase in the PCF, significant at the 1\% level throughout.

The fragmentation coefficient is where the two samples part company.
In the O*NET sample it carries the predicted negative sign in all three specifications, but the point estimates are small and none is statistically distinguishable from zero.
On O*NET occupations therefore we do not detect the workflow-level channel, but the estimates are not precise enough to rule it out either.
A modest fragmentation effect would go undetected in this sample.
On the PCF's documented sequences, however, the coefficient is negative and statistically significant in all three columns, at $-0.35$, $-0.26$ and $-0.34$, so that process groups whose AI-exposed steps are more scattered convert less of their exposure into execution.

That the channel is detected in the PCF and not in O*NET is what a test built on arrangement would lead one to expect.
The fragmentation index reads the ordering of a workflow as a whole, so it stands its best chance of detection where that ordering is recorded rather than inferred, which is the PCF's advantage over the imputed sequences and the reason we test Predictions~\#1 and~\#3 in both samples.
The trade-off noted earlier cuts the other way here.
The $525$ process groups of the PCF cover a smaller part of economic activity than O*NET's occupational structure, and they inherit their labels through a lossy match, so the channel is documented on operational business processes such as procurement, order management, after-sales service, inventory, and regulatory compliance rather than across the economy.
Neither result subsumes the other, and we read them together as saying that the arrangement channel is present where the arrangement is measured well enough to see it.

We must note that what we document in these exercises are associations that speak to the mechanism the model describes, and not causal effects.
Measuring the latter would require exogenous variation in how much of a workflow is handed to AI and/or in how that workflow is arranged, which is not what we work with.